\documentclass[12pt]{article}
\usepackage{amsmath, amssymb, amsthm}

\title{VGASP Snowman: Dissipation Proofs}
\author{Borealis S. Hedling}
\date{Version 1.0 (2026-10-08)}

\begin{document}
\maketitle

\section*{Overview}
This document provides formal proofs for dissipation behavior in the VGASP
Snowman system. Dissipation is governed by the dissipation tensor $D_{ij}$,
which encodes thermal, mechanical, and holonomy-coupled deformation. We prove:

\begin{itemize}
    \item dissipation increases monotonically under roll
    \item dissipation forbids stacking above threshold $D_{crit}$
    \item dissipation prohibits decoration on unstable surfaces
    \item dissipation couples deterministically to curvature and holonomy
\end{itemize}

All proofs rely on tensor definitions in \texttt{vgasp-snowman-math.tex}.

\section{Preliminaries}

\subsection{Dissipation Tensor}
The dissipation tensor $D_{ij}$ is defined as a symmetric, positive-semidefinite
tensor satisfying:
\[
D_{ij} = D_{ji}, \qquad D_{ij} v^i v^j \ge 0,
\]
for any vector $v^i$. Dissipation evolves according to:
\[
\dot{D}_{ij} = \Phi_{ij}(K, H),
\]
where $\Phi_{ij}$ is a tensor-valued dissipation functional.

\subsection{Critical Threshold}
Define the dissipation threshold:
\[
D_{crit} = \sup_{\Sigma} D_{ij} g^{ij},
\]
where $g^{ij}$ is the metric tensor. Operations above $D_{crit}$ are forbidden.

\section{Monotonic Dissipation Under Roll}

\begin{lemma}
Roll increases dissipation monotonically.
\end{lemma}

\begin{proof}
Roll modifies curvature $K_{ij}$ and holonomy $H_{ijk}$ along a trajectory
$\gamma(t)$. Dissipation evolves as:
\[
\frac{d}{dt} D_{ij} = \Phi_{ij}(K(\gamma(t)), H(\gamma(t))).
\]
Since $\Phi_{ij}$ is positive-semidefinite by definition,
\[
\frac{d}{dt} (D_{ij} v^i v^j) \ge 0,
\]
for any $v^i$. Therefore dissipation increases monotonically under roll.
\end{proof}

\section{Stacking Forbidden Above $D_{crit}$}

\begin{lemma}
Stacking is forbidden when $D_{ij} g^{ij} > D_{crit}$.
\end{lemma}

\begin{proof}
Stack requires curvature stability:
\[
K_{ij} n^i n^j > 0,
\]
where $n^i$ is the surface normal. Dissipation modifies curvature via:
\[
\delta K_{ij} = \Psi_{ij}(D),
\]
where $\Psi_{ij}$ is a deformation functional. If $D_{ij} g^{ij} > D_{crit}$,
then $\Psi_{ij}$ destabilizes curvature:
\[
K_{ij} n^i n^j \le 0.
\]
Thus stack violates curvature stability and is forbidden.
\end{proof}

\section{Decoration Forbidden on Dissipating Surfaces}

\begin{theorem}
Decoration cannot occur on surfaces where $\dot{D}_{ij} \neq 0$.
\end{theorem}

\begin{proof}
Decoration requires a stable surface $\Sigma$ satisfying:
\[
\mathcal{L}_{\mathcal{D}} g_{ij} = 0.
\]
If dissipation is active, then:
\[
\dot{D}_{ij} = \Phi_{ij}(K, H) \neq 0,
\]
which induces deformation:
\[
\mathcal{L}_{\mathcal{D}} g_{ij} = \Theta_{ij}(D) \neq 0.
\]
Thus decoration modifies geometry in a manner inconsistent with invariant
surfaces, violating identity constraints. Therefore decoration is forbidden.
\end{proof}

\section{Coupling of Dissipation to Curvature and Holonomy}

\begin{theorem}
Dissipation couples deterministically to curvature and holonomy.
\end{theorem}

\begin{proof}
Curvature evolves under dissipation as:
\[
\dot{K}_{ij} = A_{ij}(D),
\]
and holonomy evolves as:
\[
\dot{H}_{ijk} = B_{ijk}(D).
\]
Since $D_{ij}$ is positive-semidefinite and evolves deterministically under
$\Phi_{ij}$, both $A_{ij}$ and $B_{ijk}$ inherit deterministic evolution:
\[
\nabla_m \dot{K}_{ij} = 0, \qquad
\nabla_m \dot{H}_{ijk} = 0.
\]
Thus dissipation couples deterministically to curvature and holonomy.
\end{proof}

\section*{Conclusion}
We have shown that dissipation:

\begin{itemize}
    \item increases monotonically under roll
    \item forbids stacking above $D_{crit}$
    \item prohibits decoration on unstable surfaces
    \item couples deterministically to curvature and holonomy
\end{itemize}

These results form the foundation for JSON operator validation and
system-agnostic execution.

\end{document}
